% ============================================================
%  SentinelAPI — Plain-English Explainer (non-technical)
%  Every technical term is explained in simple words + a glossary.
%  Compile:  pdflatex nontechnical.tex
% ============================================================
\documentclass[11pt]{article}

\usepackage[margin=1in]{geometry}
\usepackage[T1]{fontenc}
\usepackage[utf8]{inputenc}
\usepackage{xcolor}
\usepackage{titlesec}
\usepackage{enumitem}
\usepackage{booktabs}
\usepackage{helvet}
\usepackage[hidelinks]{hyperref}

\renewcommand{\familydefault}{\sfdefault}
\definecolor{accent}{RGB}{37, 99, 235}
\definecolor{red}{RGB}{220, 38, 38}
\definecolor{orange}{RGB}{234, 88, 12}
\definecolor{grayline}{RGB}{229, 231, 235}

\titleformat{\section}
  {\Large\bfseries\color{accent}}{\thesection.}{0.6em}{}
\titleformat{\subsection}
  {\large\bfseries}{\thesubsection}{0.6em}{}

\newcommand{\card}[2]{%
  \noindent\fcolorbox{grayline}{white}{%
    \begin{minipage}{\dimexpr\linewidth-2\fboxsep-2\fboxrule\relax}
      \textbf{\color{accent}#1}\\[2mm]
      #2
    \end{minipage}}%
  \par\vspace{3mm}}

\title{\bfseries SentinelAPI --- Explained in Plain English\\[1mm]
       \Large A tool that finds hidden security holes in websites, automatically}
\author{AMIHACKS 2026 --- Track C (Deep-Tech)}
\date{\today}

\begin{document}

\maketitle
\thispagestyle{empty}

\begin{center}
\fcolorbox{accent}{white}{\parbox{0.92\linewidth}{%
\centering\itshape
``We built a robot that finds a website's hidden security holes automatically ---
before a hacker does.''}}\end{center}

\vspace{4mm}

\section{The Problem, in One Minute}

Apps and websites talk to servers through ``doors'' called \textbf{APIs}. Every
time you log in, check a balance, or open a profile, your app is knocking on one
of these doors. When a door is built carelessly, two things go wrong:

\begin{enumerate}[leftmargin=1.8em, itemsep=1mm]
  \item \textbf{The open door.} The server hands over \emph{someone else's} data
        when you ask --- because nobody checked \emph{who} is asking.
  \item \textbf{The oversharing door.} The server sends back \emph{too much} ---
        passwords, card numbers, private details the app never needed.
\end{enumerate}

Hackers find these mistakes every day. Companies usually notice only
\emph{after} the damage is done, because checking by hand is slow and expensive.

\section{What We Built}

\textbf{SentinelAPI} is an \emph{automatic security inspector}. You point it at
an app's doors, and it tells you exactly which ones are unsafe --- in plain
words, with a one-line proof, so anyone can act on it.

\card{In one sentence}{
It creates two imaginary users, checks whether one can see the other's secrets,
and checks whether the server is sharing more than it promised.}

\section{How It Works --- Five Simple Steps}

\begin{enumerate}[leftmargin=1.8em, itemsep=2mm]
  \item \textbf{We built a practice app} with known flaws, like a training dummy
        for the inspector to test itself on.
  \item \textbf{The inspector creates two imaginary users} --- Alice and Bob.
  \item \textbf{It pretends to be Alice and asks for Bob's stuff.}
        If the server says \emph{``sure, here's Bob's data''} --- that's a red
        flag. We call it the \emph{open-door} problem.
  \item \textbf{It compares the promise with reality.} The server's own rulebook
        says what it will share; the inspector checks what it \emph{actually}
        shares. Extra secrets --- passwords, card numbers --- are another red
        flag. We call it the \emph{oversharing} problem.
  \item \textbf{It writes a report.} Every problem gets a colour badge
        (\textcolor{red}{red} = serious, \textcolor{orange}{orange} = important)
        plus a copy-paste command that proves it.
\end{enumerate}

\section{What Makes It New}

\card{We test, we don't guess}{
Most scanners throw random junk at a server and hope something breaks. We do
something smarter: we create \textbf{two real users} and check whether one can
see the other's data. That is a \emph{real test}, not a guess.}

\card{We use the server's own rulebook}{
Every server has a ``rulebook'' (a file describing what it should share). We
compare what the server \emph{promised} against what it actually does. The
difference between promise and reality is where the leak is.}

\card{Every finding comes with proof}{
No vague warnings. Each problem includes a copy-paste command anyone can run to
see the leak with their own eyes --- so there are \emph{no false alarms}.}

\card{Runs anywhere, costs nothing}{
The whole thing works offline with free, open tools --- no paid services, no
cloud, no internet needed for the demo.}

\section{The Tech Stack, in Plain Terms}

\begin{center}
\begin{tabular}{ll}
\toprule
\textbf{Tool} & \textbf{What it does for us} \\
\midrule
Python + FastAPI & The practice app, and the ``brain'' of the inspector \\
Go & A second, faster copy of the brain (single file, no setup) \\
OpenAPI & The server's rulebook the inspector reads \\
HTML + CSS + Three.js & The clean report page with colour badges and a 3D background \\
\bottomrule
\end{tabular}
\end{center}

\section{The Results}

We tested the inspector on \textbf{two practice apps} with 24 hidden problems in
total. It found \textbf{all 24}, and raised \textbf{zero} false alarms. In
measuring terms: \textbf{100\% precision, 100\% recall}. In plain words: it
catches everything real and never cries wolf.

\section{Jargon, Translated}

\noindent The technical words, explained simply:

\begin{center}
\small
\begin{tabular}{ll}
\toprule
\textbf{Term} & \textbf{What it actually means} \\
\midrule
API & The ``door'' an app uses to ask a server for data \\
BOLA / IDOR & The ``open door'' bug --- seeing other people's data \\
OpenAPI & The server's written ``rulebook'' of what it should share \\
Differential testing & Creating two users and comparing what each can see \\
Broken authentication & The door works with \emph{no} key at all \\
Privilege escalation & A normal user getting \emph{admin} powers they shouldn't have \\
Mass assignment & The server accepting extra fields it was never asked for \\
SSRF guard & A safety lock so our tool can't be abused to attack other servers \\
Precision / Recall & How often it's right vs.\ how many real bugs it misses \\
\bottomrule
\end{tabular}
\end{center}

\vspace{4mm}
\begin{center}
\fcolorbox{grayline}{white}{\parbox{0.92\linewidth}{%
\centering\large\bfseries
The result: a tool that finds security holes automatically, explains them in
plain words, and proves them --- before a hacker does.}}\end{center}

\end{document}
